\begin{enumerate}
\item Nejprve si stáhněte soubor přiložený k~diplomové práci z~OneDrive, který obsahuje jak samotný projekt, tak také pluginy, jež je nutné přidat do Unreal Enginu.
\item Pokud nemáte nainstalovaný \uv{Epic Games Store}, stáhněte si jej ze stránky \texttt{https://store.epicgames.com/}. V~pravém horním rohu klikněte na tlačítko \uv{Download} a~spusťte instalátor, čímž nainstalujete \uv{Epic Games Launcher}.
\item Po spuštění aplikace si vytvořte účet (pokud jej ještě nemáte), buď přímo v \uv{Epic Games Store}, nebo prostřednictvím \uv{Epic Games Launcher}, a~přihlaste se.
\item Po přihlášení klikněte v~levém menu na položku \uv{Unreal Engine}.
\item V~horní části aplikace se zobrazí nové možnosti — vyberte záložku \uv{Library}.
\item V~sekci \uv{Engine Versions} klikněte na tlačítko \uv{+}. Objeví se nabídka výběru verze enginu. Zvolte verzi \texttt{5.2.1} a~spusťte instalaci. Umístění instalace můžete ponechat výchozí, ale zapamatujte si ho — bude potřeba později.
\item Po dokončení instalace přejděte do složky s~nainstalovaným enginem, která se bude jmenovat \texttt{UE\_5.2}. Do této složky vložte podsložku \texttt{Engine} ze složky \uv{Pluginy}, která je umístěna vedle složky projektu \uv{RoboHell}. Tímto krokem zajistíte správné nainstalování potřebných pluginů.
\item Projekt spusťte dvojklikem na soubor \uv{RoboHell.uproject} ve složce projektu. První spuštění může trvat déle, protože probíhá příprava enginu a~kompilace projektu.
\item Po načtení projektu jej zatím nespouštějte pomocí zeleného trojúhelníku! Před prvním spuštěním je třeba provést ještě několik kroků.
\item V~levém horním rohu klikněte na \uv{Edit} a~v~rozbaleném menu vyberte \uv{Editor Preferences}.
\item V~zobrazeném nastavení sjeďte v~levém panelu dolů k~sekci \uv{Plugins} a~rozklikněte \uv{NevarokML}.
\item Klikněte na tlačítko \uv{Refresh Details}. Vizuálně se nic nezmění, ale tímto krokem se plugin správně načte v~editoru.
\item Tímto by mělo být vše připraveno. Projekt by nyní měl být plně funkční. Doporučuji však přečíst také kapitolu \uv{Jak opravit problém s~projektem}, která řeší častou chybu při spouštění.
\end{enumerate}